\chapter*{Prefaci}
\addcontentsline{toc}{chapter}{Prefaci}

El present document constitueix la memòria de proposta del projecte pràctic de l'assignatura de Tecnologia Multimèdia, cursada al Grau d'Enginyeria Informàtica de la Universitat de les Illes Balears (UIB). El projecte rep el nom de \textit{Dona'm Bauxa}, expressió mallorquina que convida a la festa i al gaudi cultural, i proposa el disseny conceptual i l'anàlisi de requisits d'una plataforma web dedicada a la difusió i centralització d'informació sobre música en directe, artistes locals i esdeveniments culturals de l'illa de Mallorca.

La motivació per emprendre aquest projecte neix de l'observació directa d'una mancança real en l'ecosistema digital mallorquí: malgrat la riquesa i la vitalitat de l'escena musical de l'illa, no existeix cap recurs en línia especialitzat que aglutini, de manera ordenada i accessible, la informació sobre concerts, grups emergents, festivals i actes de cultura popular. Aquesta dispersió informativa perjudica tant el públic que cerca on gaudir de música en directe com els propis artistes i grups que lluiten per donar-se a conèixer.

Amb \textit{Dona'm Bauxa}, es proposa una solució concreta i vigorosa a aquesta problemàtica, apostant per la proximitat, el talent local i l'ús de tecnologies web modernes com a eines de cohesió cultural. Al llarg d'aquest document s'exposa de manera detallada la identificació del grup de treball, la justificació de la temàtica escollida, l'anàlisi funcional del sistema, allò que farà la plataforma, i l'anàlisi no funcional, és a dir, com haurà de funcionar per satisfer les expectatives dels seus usuaris.

Volem expressar el nostre agraïment al professorat de l'assignatura per la seva guia i pels recursos posats a disposició de l'alumnat, i a tots els músics i artistes mallorquins que, amb la seva feina, fan de la nostra illa un lloc on la bauxa mai s'acaba.

\begin{flushright}
{\makeatletter\itshape
    Josep Ferriol Font \& Dylan Luigi Canning Garcia \\
    Palma de Mallorca, febrer del 2026
\makeatother}
\end{flushright}
